\section{A Streaming Garbled-Circuit Factory}
\label{sec:ongoing}

The finite construction samples one bounded garbling.  For ongoing
computation we now describe a more specific target than a general universal
streamer: an \emph{oblivious garbled-circuit sampler} (OGCS) for one fixed
public transition circuit
\[
  F(s_e,a_e)=(s_{e+1},y_e),
\]
where $s_e$ is carried state, $a_e$ is fresh input, and $y_e$ is public.  One
application of $F$ is a \emph{stretch}.  The desired public object has a finite
description and accepts arbitrarily long finite histories; in particular, its
interface has no setup-time horizon.

This section gives a candidate construction and a precise interface against
which a proof can be written.  It is not yet a theorem.  The cited
general-purpose ordinary-iO-based succinct TM-obfuscation results do not
immediately provide the horizon-free interface, and the final online
random-oracle programming step also remains open in this first version.

\subsection{The prefix interface and its schedule}

Setup publishes a finite parameter $U_{\mathsf{ogcs}}$ fixing $F$ and
containing a fresh identifier $\mathsf{uid}$.  The public algorithm is
evaluated on a complete finite input history:
\[
\begin{split}
  \mathsf{OGCS}(U_{\mathsf{ogcs}};I,
     &(M_0,\Sigma_0,\mathbf W_0),\ldots,\
       (M_e,\Sigma_e,\mathbf W_e))
       \longrightarrow D_e .
  \label{eq:ogcs-interface}
\end{split}
\]
The physical object is copyable and stateless; its hidden \emph{logical}
state is reconstructed by replaying the input history.  On every accepting
history its outputs are exactly prefix preserving,
\[
  D_{e-1}\preceq D_e,
\]
and repeated evaluation on the same history gives the same bit string.  A
failed check enters an absorbing reject state and emits no further extension.
The certificate is an acceptance witness, not stream input: two histories are
called command-equivalent when they have the same $I$ and the same sequence of
$(M_e,\mathbf W_e)$ values, even if they use different valid encodings of
$\Sigma_e$.  The program normalizes away certificate bytes after verification.
A different initialization or non-equivalent accepted command history is a
fork and must not expose the hidden state of the canonical history.

The initialization input $I$ produces a preparation object
$\mathsf{Prep}_0$ but no garbled table or application output.  For a fixed
public initial state $s_0$, it also publishes exactly one active initial-state
label per state wire as part of $\mathsf{Prep}_0$.  The first certificate key
$\mathsf{VK}_0$ is likewise a component of $\mathsf{Prep}_0$.  In general,
$\mathsf{Prep}_e$ is public before the input
slots of stretch $e$ open.  It contains
\begin{itemize}
  \item encrypted activation tokens and verifiable label shares, together
    with public token-validity tags, for the fresh external-input wires of
    stretch $e$; and
  \item encrypted branch secrets and public verification material for a
    fresh one-time certificate used to close that stretch.
\end{itemize}
The hidden plaintext material underlying $\mathsf{Prep}_e$ was fixed by the
preceding accepted extension, or by $I$ when $e=0$.  For notation set
$D_{-1}:=\mathsf{Prep}_0$.

The order within a stretch is important.
\begin{enumerate}
  \item Input providers post activation tokens.  At a deterministic ledger
    cut every slot is closed and the selected input committees reconstruct one
    active external-input label per wire.  Missing or malformed slots use the
    publicly detectable random-default rule of
    \cref{sec:basic-construction}.
  \item The registration window for stretch $e+1$ closes.  The canonical
    ledger cut determines one ordered vector $\mathsf{PKBatch}_{e+1}$.  If the
    required number of fresh keys is unavailable, the extension waits.
  \item Everyone computes the finite vector $\mathbf W_e$ of random-oracle
    values specified below.  This vector may be transported by any party and
    need not be stored on the ledger; it has no authority of its own.
  \item The certificate roles derive one canonical checkpoint $M_e$ from the
    closed ledger view and reveal a certificate $\Sigma_e$ for it.
  \item The accepted triple $(M_e,\Sigma_e,\mathbf W_e)$ appends
    \[
      \GC_e,\qquad \mathsf{Dec}_e,\qquad \mathsf{Prep}_{e+1}
      \label{eq:stretch-output}
    \]
    to the cumulative OGCS output.
  \item Everyone evaluates $\GC_e$.  This gives the public output $y_e$ and
    one active label on every next-state wire.
\end{enumerate}
Publishing the garbled tables only after input closure lets the straight-line
simulator obtain the public ideal output before it must fix those tables.  A
variant publishing the complete garbling earlier would require a stronger
online or adaptively simulatable garbling notion.

\subsection{A one-time common-input certificate from silence}
\label{sec:yamso-certificate}

The ledger has no signing key.  Persistence instead comes from fresh YAMSO
roles that certify a value they compute from their common ledger view.

Let
$\mathsf{Code}:\{0,1\}^{\kappa}\rightarrow\{0,1\}^{n}$ be a binary code of
relative distance $\delta$, let $h_{\mathsf{sig}}$ be collision resistant,
and let $f$ be one way.  For certificate epoch $e$, the hidden generator
samples independently
\[
  A^e_{j,b}\sample\{0,1\}^{\lambda},\qquad
  Y^e_{j,b}=f(\mathsf{sid},e,j,b,A^e_{j,b})
\]
for $(j,b)\in[n]\times\{0,1\}$.  The public verification material is
$\mathsf{VK}_e=(Y^e_{j,b})_{j,b}$.  Each preimage $A^e_{j,b}$ is encrypted to
a different fresh one-use computation role $P^e_{j,b}$.

At the fixed certification cut, every honest role independently computes the
same canonical message $M_e$ and the codeword
\[
  c_e=\mathsf{Code}(h_{\mathsf{sig}}(M_e)).
\]
The protected algorithm of role $P^e_{j,b}$ returns its secret key precisely
when $c_e[j]=b$.  Observers check $\mathsf{KeyMatch}$, decrypt the canonical
role packet, and treat the recovered $A^e_{j,b}$ as the opening.  A
certificate contains valid openings in at least $\tau n$
coordinates matching the claimed codeword.  Corrupt matching roles may
withhold, which is why verification does not require every coordinate.

Set $\rho=p$, the computation-role corruption probability of the model, and
choose constant margins
$\gamma_{\mathsf{live}},\gamma_{\mathsf{forge}}>0$ such that
\[
  1-\delta(1-\rho)+\gamma_{\mathsf{forge}}
  <\tau<
  1-\rho-\gamma_{\mathsf{live}}.
  \label{eq:certificate-window}
\]
The right inequality gives liveness.  For a different codeword, every
disagreeing coordinate needs the opposite opening, which is available without
inverting $f$ only when that opposite role is corrupt; the left inequality
prevents enough such openings.  The interval is possible when
$\delta>\rho/(1-\rho)$, and hence permits any $\delta>1/3$ at
$\rho=1/4$.  The canonical digest may depend on corrupt-role information
visible before its cut.  We therefore require the good-assignment event to
hold simultaneously for all possible canonical codewords and all ordered
pairs of distinct codewords.  A conservative sufficient condition is
\[
  2^{\kappa}\exp(-\Omega(\gamma_{\mathsf{live}}^2 n))
  +2^{2\kappa}\exp(-\Omega(
       \gamma_{\mathsf{forge}}^2\delta n))
\]
to be negligible.  The exact constants and concentration proof are part of
the certificate lemma; they are not claimed proved here.

Speaking exposes the complete state of $P^e_{j,c_e[j]}$, but its branch secret
has just become public.  The role containing $A^e_{j,1-c_e[j]}$ remains
silent forever.  This is precisely where the at-most-once rule creates
persistent authenticity without a persistent signing key.

The distance-amplified Lamport idea is not new.  Katz and Vaikuntanathan use a
high-distance encoding in a leakage-resilient one-time signature
\cite{KatzVaikuntanathan2009LeakageSignatures}.  We reuse the distance
intuition, while the distributed role assignment, partial liveness threshold,
post-speech full exposure, and common-input property require a separate
argument.  Kelsey, Lang, and Lucks study distributed stateful hash-based
signing and the disjoint-honest-trustees one-time-key-reuse problem in the
absence of broadcast \cite{KelseyLangLucks2025DistributedHBS}.

The certificate is not a signing service for adversarially supplied messages.
Honest roles derive $M_e$ from one already closed ledger cut.  The common
ledger prevents them from being shown different inputs; the certificate makes
that public truth usable by an obfuscated program that cannot read the ledger
itself.

\subsection{Keyless bootstrap and certified commands}

The initialization string has the form
\[
  I=(\mathsf{Init},\mathsf{uid},\mathsf{sid},
       \mathsf{initDigest},\mathsf{PKBatch}_0).
\]
The program requires the supplied $\mathsf{uid}$ to equal the one embedded in
$U_{\mathsf{ogcs}}$.  The canonical $\mathsf{initDigest}$ commits to the
initial ledger cut, public initial state, role allocation, and all relevant
parameters.  The hidden program absorbs the complete encoding of $I$ before
deriving its root and emitting $\mathsf{Prep}_0$, which includes
$\mathsf{VK}_0$ and the active labels for the fixed $s_0$.  Thus initialization
is not literally outputless: it emits branch-local encrypted preparation
material and public labels for a public state, but no garbled table or
application-derived output.

Anyone may run the copyable program on $I'\neq I$.  Such a call may even use
only adversarial public keys and hence create an adversarially controlled
certificate chain.  This is harmless only because the different
initialization is absorbed before any branch material is derived: every
token, label, ciphertext, and later certificate key on that fork lies under a
pseudorandomly independent logical root.  Later exposure of a canonical role
key may decrypt a packet addressed to the same public key on such a fork, but
the plaintext contains only fork material.  The common ledger selects the
live root; no certifying-ledger secret key is needed.

For a public partition of microstep indices into finite intervals $J_e$, let
\[
  W_i=H(\mathsf{ogcs},\mathsf{uid},\mathsf{sid},i),
  \qquad i\in J_e,
\]
write $\mathbf W_e=(W_i)_{i\in J_e}$, and set
$\omega_e=\mathsf{Hash}(J_e,\mathbf W_e)$.  The obfuscated program has no
random-oracle access, so these values are computed outside it.  Honest
certificate roles compute the same vector themselves.  The vector supplied
with the OGCS call is nonauthoritative: the program checks its digest against
the certified $\omega_e$.

Define the full pre-oracle descriptor and its digest by
\[
  Q_e=(\mathsf{uid},\mathsf{sid},e,
    \mathsf{Hash}(D_{e-1}),\mathsf{Hash}(\mathsf{Prep}_e),
    \mathsf{inputCheckpoint}_e,\mathsf{PKBatch}_{e+1}),
  \qquad
  \chi_e=\mathsf{Hash}(Q_e),
\]
and let the canonical certificate message be
\[
  M_e=(Q_e,J_e,\omega_e).
  \label{eq:checkpoint}
\]
The next-key registration cutoff occurs before this message is certified.
Exact closed-ledger positions, role allocation, and all domain separators are
part of $Q_e$'s canonical encoding.  Thus the program receives, checks, and
can use the actual input checkpoint and recipient-key batch rather than only
their hash image.  It checks the predecessor, epoch, syntax, oracle-vector
digest, and certificate under the verification material in its current
logical frontier.  Any failure enters the absorbing reject state.

The next verification key is an output inside $\mathsf{Prep}_{e+1}$ and hence
is authenticated by the current accepted extension.  It is not included as an
independent input to the same command, avoiding a circularity.  Different
valid subsets or orderings of openings for the same $M_e$ induce the same
extension: all derivations use $M_e$, never the certificate bytes.

\subsection{The hidden logical stream}

Let $\mathsf{InMat}_e$ contain the external-input label pairs, activation
tokens, semantic permutations, shares, and role-packet plaintexts underlying
$\mathsf{Prep}_e$, and let $\mathsf{CertMat}_e$ contain its certificate
preimages and packet state.  At the beginning of stretch $e$, the logical
frontier is
\[
  \Phi_e=\left(
    X_e,
    \mathsf{StatePairs}_e,
    \mathsf{InMat}_e,
    \mathsf{CertMat}_e,
    \mathsf{VK}_e
  \right),
  \label{eq:logical-frontier}
\]
After the public cuts are fixed, let $\overline\Phi_e$ be the frontier
$\Phi_e$ with its seed coordinate replaced by
\[
  \overline X_e=
    \mathsf{KDF}(X_e,(\mathsf{publicCut},\chi_e)).
\]
Write $\Phi_{e,i}$ for the corresponding intermediate frontier after
microstep $i$, with $\Phi_{e,\ell_e-1}=\overline\Phi_e$ and
$\Phi_{e,r_e}=\Phi_{e+1}$ on a normally generated stretch.  Here
$\mathsf{StatePairs}_e=
 \{(K^{\mathsf{st}}_{e,w,0},K^{\mathsf{st}}_{e,w,1})\}_w$.
The active member of every state pair is public from evaluation of the
preceding stretch.  For $e=0$, $\mathsf{Prep}_0$ publishes the member encoding
the fixed public initial state; the sibling remains hidden.

Enumerate the finite objects generated for the next suffix by the consecutive
indices $J_e=[\ell_e,r_e]$.  Set
\[
  X_{e,\ell_e-1}=\overline X_e,
  \qquad
  X_{e,i}=\mathsf{KDF}(X_{e,i-1},W_i)
  \quad(i\in J_e),
  \qquad
  X_{e+1}=X_{e,r_e}.
  \label{eq:kdf-chain}
\]
These equations specify the normal branch.  At a programmed microstep, a
valid capsule replaces both the emitted fragment and the complete successor
microfrontier $\Phi_{e,i}$; subsequent generation continues from that
replacement frontier.
Domain-separated puncturable-PRF evaluations from these states derive the
output-state label pairs and garbling coins for $\GC_e$, followed by
$\mathsf{InMat}_{e+1}$, $\mathsf{CertMat}_{e+1}$, verification material,
verifiable shares, and role-state encryption coins for
$\mathsf{Prep}_{e+1}$.  The current $\mathsf{InMat}_e$ is not regenerated:
it is the material already committed before the input window.  The complete
pre-oracle cut is therefore absorbed before the first microstep.

Readout for the newly prepared input roles is bound to the exact ordered
recipient-key vector.  For example, its coins have the form
\[
 r^{\mathsf{read}}_{e+1,w,s,j}
 =\mathsf{PPRF}_{K_{\mathsf{read}}}
   (X_{e,i},\chi_e,e+1,w,s,j,
    pk_{e+1,w,s,j}).
  \label{eq:pk-bound-readout}
\]
Consequently, a canonical packet is tied to its intended recipient key.
Packets produced by reset or fake-initialization calls under the same public
key are not claimed never to decrypt; rather, prefix isolation must ensure
that their plaintexts contain only independently rooted fork material.

The garbling $\GC_e$ uses $\mathsf{StatePairs}_e$ and the external-input pairs
in $\mathsf{InMat}_e$.  Its output-state labels are literally the pairs in
$\mathsf{StatePairs}_{e+1}$ stored in $\Phi_{e+1}$.  Evaluation makes only the
active member of each pair public.  No machine links an old label to a new
one, and no committee is needed for carried state.

Changing $I$, an accepted certified checkpoint, or any earlier $W_i$ changes
the logical derivation for the suffix after that point.  On the live protocol
path, the certificate prevents such a change after a shared state label
exists.  Calls on another initialization root remain independently evaluable
but receive no honest input material belonging to the live root.  In a
programmed proof hybrid, a sequence of transition capsules may prescribe each
affected public fragment together with its matching successor microfrontier;
no capsule may program only a table while leaving an inconsistent hidden
successor.

\subsection{Candidate realization by delayed programming}

For a declared input-length bound $L$, a finite version of the public object
may be based on succinct iO for bounded-input Turing machines.  Cryptographic
iterators and succinct TM-iO provide relevant techniques for compactly
representing long bounded computations
\cite{KoppulaLewkoWaters2014IOTM,AnanthJainSahai2017TMIO}.  They do not by
themselves instantiate the horizon-free interface above, whose replay input
has unbounded length; that requires an additional primitive or a
reformulation of the interface.  We leave this as a separate realization gap
rather than hide $L$ in the environment's running-time polynomial.

For a bounded prefix, we propose to adapt the delayed-backdoor programming of
Hofheinz et al.'s universal sampler
\cite{HofheinzEtAl2014UniversalSamplers}.  Its hidden-trigger paradigm, and
the later selectively secure pseudorandom puncturable deterministic encryption
(PPDE) of Koppula, Poelstra, and Waters
\cite{KoppulaPoelstraWaters2017FastSamplers}, let an apparently random string
carry a programmed bounded object.  In our proposed adaptation, a random
$W_i$ selects the ordinary PRF-generated transition.  In a proof, the oracle
answer at one punctured point is changed to encode a prescribed transition
capsule.

More explicitly, let ($\mathsf{BEnc},\mathsf{BDec}$) denote the desired
puncturable carrier.  Its ciphertexts are pseudorandom, while a uniform string
decrypts to $\bot$ except with negligible probability.  The same hidden
microstep program is used in both the real execution and the proof:
\[
  C_i\leftarrow\mathsf{BDec}_{K_{\mathsf{bd}}}(W_i),
  \qquad
  \mathsf{Step}_i=
  \begin{cases}
    C_i,&\text{if }\mathsf{ValidCap}(C_i,\Phi_{e,i-1},\chi_e)=1,\\
    \mathsf{NormalGen}(X_{e,i-1},W_i,\chi_e,i),&\text{otherwise.}
  \end{cases}
  \label{eq:normal-programmed-step}
\]
On honest random-oracle answers, the normal branch is taken except with
negligible probability.  In a selective bounded-prefix hybrid in which
$C_i$ is fixed before carrier-key setup, the intended proof first replaces
the embedded carrier key by a key punctured at $C_i$ and hardcodes the
prescribed behavior at the exceptional ciphertext.  It separately punctures
the normal-generation keys at the same logical point.  Koppula--Poelstra--
Waters PPDE pseudorandomness can then hide the swap from uniform $W_i$ to
$\mathsf{BEnc}_{K_{\mathsf{bd}}}(C_i)$.  Choosing $C_i$ online after
carrier-key setup requires an adaptive carrier argument and remains open.

The programmed branch bypasses the normal KDF and readout derivations for
that microstep and takes the complete successor microfrontier from its
capsule.  Thus, in the selective experiment, the capsules can be sampled from
$\chi_e$, the closed recipient keys, and the prescribed output before $W_i$
or $M_e$ is fixed.  This is a specification of the carrier property we need,
not a claim that the cited PPDE theorem already supplies the complete stateful
interface.

For our use each capsule must contain the \emph{complete} affected logical
microtransition: its public garbled fragment and every hidden successor value
on which later fragments depend.  Programming only one public gate while
leaving an inconsistent successor $X_{e,i}$ would not preserve the stream.
The intended hybrid punctures the KDF/PRF path at one canonical microstep,
installs this complete capsule, leaves every sibling subtree pseudorandom,
and then repeats from left to right.

Two further proof obligations remain visible.  First, computational
unforgeability of the YAMSO certificate does not make two obfuscated programs
functionally identical on every mathematical input.  In their binding setup,
Fernando et al.'s somewhere statistically unforgeable signatures make
alternate valid signatures at a selected round and message statistically
nonexistent, except with negligible probability over setup.  They use this
property to establish functional equivalence of adjacent obfuscated gate
programs before applying ordinary iO \cite{FernandoEtAl2023SSU}.  This is a
candidate technique here, not an instantiation we may assume for free.

Second, the points
$H(\mathsf{ogcs},\mathsf{uid},\mathsf{sid},i)$ are predictable after
$U_{\mathsf{ogcs}}$ is public.  A selective-capsule experiment can fix the
target $\mathsf{sid}$, finite index set $\bigcup_eJ_e$, canonical schedule,
and complete capsule sequence before publication, program all corresponding
points, and then rely on the fresh $\mathsf{uid}$ to rule out an earlier
query.  This is structural evidence for the bounded carrier hybrid, not a
real--ideal simulation: in the protocol the honest inputs and outputs are not
known before setup.  It also does not give one online simulator for the
horizon-free execution.

The intended next pass prepares a fresh unpredictable nonce one stretch
ahead.  In that proof the simulator first waits for the input and next-key
cuts, receives $y_e$ privately, and uses $Z.\mathsf{Complete}$ and
$Z.\mathsf{Prepare}$ to form the microstep capsules from $\chi_e$, the
recipient keys, and $y_e$.  It then programs the nonce-dependent oracle
points, reveals and recomputes $\mathbf W_e$, forms $M_e$, and only then
releases the honest openings that take $\Sigma_e$ across its acceptance
threshold.  Thus no one can invoke the accepted canonical extension before
all of its capsules are installed.
We do not repair either issue by giving the environment an explicit time bound.

\subsection{What the candidate establishes and what it does not}

The construction isolates a plausible path to indefinitely extensible
garbling.  Adjacent stretches share boundary labels by construction; new
public keys enter only through a certified cutoff; external oracle values are
bound by a certified digest; and the only online secret holders that speak
reveal branch-local state already made public.  A fake initialization can
generate a whole fake computation, but it forks before any live input or
state material.

The missing statements are now explicit: the one-time certificate lemma, a
two-stage finite garbling/readout simulation with prescribed boundary labels
and full speaker exposure, a horizon-free public-program realization, the
stateful delayed-programming lemma, exact iO-compatible authentication, and
late online oracle programming.  These are organized as whole-process
replacement lemmas in \cref{sec:proof-outline}.  A general universal-streamer
abstraction can be revisited after this fixed-function OGCS is understood;
the present construction does not claim one.
